\documentclass[../../main.tex]{subfiles}
\begin{document}
\chapter{Notation and reusable patterns}
\label{app:notation}

\section{Frequently used macros}
Most commands in \texttt{macro.sty} are intended for math mode. The table
below gives a small selection; the package source contains additional
notation for algebra, geometry, topology, number theory, and set theory.

\begin{fullwidth}
\begin{center}
\renewcommand{\arraystretch}{1.3}
\begin{tabular}{ll}
  \toprule
  Source & Result \\
  \midrule
  \verb|\NN, \ZZ, \QQ, \RR, \CC| & $\NN,\ZZ,\QQ,\RR,\CC$ \\
  \verb|\abs{x}, \norm{v}| & $\abs{x},\norm{v}$ \\
  \verb|\floor{x}, \ceiling{x}| & $\floor{x},\ceiling{x}$ \\
  \verb|\cbrt{x}, \half, \eps| & $\cbrt{x},\half,\eps$ \\
  \verb|\Hom(G,H), \Aut(G)| & $\Hom(G,H),\Aut(G)$ \\
  \verb|\Ker f, \img f, \rank A| & $\Ker f,\img f,\rank A$ \\
  \verb|\Zc{n}, \Zm{n}| & $\Zc{n},\Zm{n}$ \\
  \verb|\injto, \surjto, \taking{f}| & $\injto,\surjto,\taking{f}$ \\
  \verb|\fpartial{f}{x}| & $\fpartial{f}{x}$ \\
  \verb|\arc{AB}, \seg{AB}, \ray{AB}| & $\arc{AB},\seg{AB},\ray{AB}$ \\
  \bottomrule
\end{tabular}
\end{center}
\end{fullwidth}

\begin{notation}
Here $\NN=\{0,1,2,\ldots\}$. The macro \verb|\NN| only prints a symbol;
you must state whether your own convention includes zero.
\end{notation}

\section{Project-specific notation}
The file \texttt{macros/macros.tex} defines two additional commands used
only in this book. The command \verb|\angles{v,w}| prints
$\angles{v,w}$, and \verb|\dd| prints an upright differential, as in
$\int_0^1 x\,\dd x$.

Define extra commands in that file with \verb|\newcommand|. If a name is
already defined, choose a more specific name instead of silently replacing
an existing command. In particular, \texttt{macro.sty} already changes the
meanings of \verb|\AA|, \verb|\th|, \verb|\Re|, and \verb|\Im|.

\clearpage
\section{A writing checklist}
Before sharing a chapter:
\begin{enumerate}[label=\arabic*.]
  \item State the domain of each variable and any conventions you use.
  \item Give definitions before using their terminology in results.
  \item Label equations, theorems, figures, and tables that you reference.
  \item Give figures informative captions and cite external sources.
  \item Build the whole book to resolve references between chapters.
  \item Check the log for missing citations and undefined references, then
        inspect the final PDF for layout problems.
\end{enumerate}

\begin{note}
A chapter compiled on its own has its own counters and auxiliary files.
Use the complete book when checking final chapter numbers, cross-references,
the table of contents, and the bibliography.
\end{note}
\end{document}
